\documentclass[11pt,a4paper]{article}
\usepackage[T1]{fontenc}
\usepackage{amsmath,amssymb,geometry,booktabs,parskip,xcolor,hyperref}
\geometry{margin=2.3cm}
\hypersetup{colorlinks=true,linkcolor=blue!60!black,urlcolor=blue!60!black}
\title{\textbf{Work plan}\\\large Quantum resources for multi-asset optimal execution}
\author{Research route and notes for the paper}
\date{}
\begin{document}
\maketitle

\section*{Purpose and current position}

Estimate the logical qubits and T gates needed to solve an execution
problem that balances the cost of selling quickly against the price risk
of holding inventory. The route is
\[
 \text{execution HJB}\longrightarrow\text{Cole--Hopf linearisation}
 \longrightarrow\text{Schr\"odingerisation}\longrightarrow\text{resource count}.
\]
The two-asset execution problem has a classical Riccati solution, which
provides the reference answer. The contribution is a resource estimate;
no quantum speedup is claimed.

\textbf{Current: Step 3 verification. Next: complete Step 4.}
Step 1 is not needed. Step 2 starts from the published Eq.~16 with $R=0$.
No numerical PDE validation or complete quantum resource count is available.

\texttt{WorkPlan.tex} holds derivations, sources and checks for the paper;
\texttt{Summary.tex} holds equations and one-line statuses.

\subsection*{Model and notation}

$\mathbf x\in\mathbb R^M$ is inventory, $\mathbf X=\mathbf x(0)$,
$\mathbf v=-\dot{\mathbf x}$ is selling rate, and $T$ is the horizon.
$H=\operatorname{diag}(\eta_j)\succ0$ is the impact matrix,
$C\succeq0$ is price covariance, and $\lambda\ge0$ penalizes risk.
Coefficients are constant, controls are unrestricted signed rates, and no
spread cost or participation limit is included. Sales-only feasibility is
checked for the chosen starting portfolio in Step 3.

A finite terminal penalty is $\tfrac12\mathbf x(T)^\top\Phi\mathbf x(T)$,
$\Phi\succeq0$. Exact liquidation instead requires $\mathbf x(T)=0$.
These are different boundary-value problems. The diffusion parameter $\nu$
introduced in Step 4 acts in inventory space and is not price volatility.
The impact matrix $H$ and the differential operator $\hat H$ are distinct.

\section{Step 2: specialize the 2022 paper's Eq.~16}
\label{sec:stage2}

\subsection*{Source and reduction}

Use Bergault, Drissi and Gu\'eant \cite{bergault2022}, Eq.~16 with terminal
condition Eq.~17, Remark~2 and Appendix~A. Equation numbers refer to
arXiv:2103.13773v4, 22 March 2022, p.~7.

In the paper's notation, Eq.~16 is
\begin{equation}
\begin{aligned}
0={}&\partial_t\theta
 +\tfrac14(\nabla_q\theta)^\top\eta^{-1}\nabla_q\theta
 +\tfrac12\operatorname{tr}(\Sigma D^2_{SS}\theta)\\
 &-\tfrac\gamma2(q+\nabla_S\theta)^\top\Sigma(q+\nabla_S\theta)
 +(\bar S-S)^\top R^\top(\nabla_S\theta+q).
\end{aligned}
\label{eq:bdg16}
\end{equation}
Set $R=0$. As stated in Remark~2 and Appendix~A of the paper, its quadratic
solution then has no price dependence: $\theta(t,q,S)=q^\top A(t)q$.
Thus $\nabla_S\theta=0$ and $D^2_{SS}\theta=0$, and Eq.~16 becomes
\[
0=\partial_t\theta+\tfrac14(\nabla_q\theta)^\top\eta^{-1}\nabla_q\theta
       -\tfrac\gamma2q^\top\Sigma q.
\]
Translate to this project's cost-minimization convention using
\[
 q=\mathbf x,\qquad \theta=-\mathcal V,\qquad
 \eta=H,\quad\Sigma=C,\quad\gamma=2\lambda,\quad\Gamma=\Phi/2.
\]
Multiplying the reduced equation by $-1$ gives Eq.~\eqref{eq:pde} below.
The paper uses $\dot q=v_{\rm paper}$, whereas we use
$\dot{\mathbf x}=-\mathbf v$. Hence
$v_{\rm paper}^*=\eta^{-1}\nabla_q\theta/2=-\mathbf v^*$, which gives
the feedback in Eq.~\eqref{eq:pde}. Its terminal condition
$\theta(T,q,S)=-q^\top\Gamma q$ becomes
$\mathcal V(\mathbf x,T)=\tfrac12\mathbf x^\top\Phi\mathbf x$.
No discrete-to-continuum derivation is needed for this route.

\subsection*{Classical reference solution}

The corresponding reduced cost-to-go is
\begin{equation}
 \mathcal V_\Phi(\mathbf x,t)=\inf_{\mathbf v(\cdot)}
 \left\{\int_t^T(\mathbf v^\top H\mathbf v+
 \lambda\mathbf x^\top C\mathbf x)\,dr
 +\tfrac12\mathbf x(T)^\top\Phi\mathbf x(T)\right\},
 \qquad\dot{\mathbf x}=-\mathbf v.
 \label{eq:cost}
\end{equation}
For deterministic schedules, Gaussian cost gives a certainty-equivalent risk
term $\gamma\operatorname{Var}/2$, hence $\gamma=2\lambda$.
The specialized HJB and feedback are
\begin{equation}
 \boxed{\mathcal V_t-\tfrac14(\nabla\mathcal V)^\top H^{-1}\nabla\mathcal V
 +\lambda\mathbf x^\top C\mathbf x=0,\qquad
 \mathbf v^*=\tfrac12H^{-1}\nabla\mathcal V.}
 \label{eq:pde}
\end{equation}
For $g=\nabla\mathcal V$, the identity
\[
 \mathbf v^\top H\mathbf v-\mathbf v^\top g
 =(\mathbf v-\tfrac12H^{-1}g)^\top H(\mathbf v-\tfrac12H^{-1}g)
 -\tfrac14g^\top H^{-1}g
\]
shows that the feedback uniquely minimizes the Hamiltonian.
Substitute $\mathcal V=\tfrac12\mathbf x^\top P(t)\mathbf x+c(t)$ to obtain
\begin{equation}
 \dot P=\tfrac12PH^{-1}P-2\lambda C,\qquad P(T)=\Phi,
 \qquad\dot c=0,\quad c(T)=0,
 \qquad\mathbf v^*=\tfrac12H^{-1}P\mathbf x.
 \label{eq:ric}
\end{equation}
Thus $c=0$ and the optimized initial cost is
$\mathcal V_\Phi(\mathbf X,0)=\tfrac12\mathbf X^\top P(0)\mathbf X$.
Along an admissible path,
\[
 J_\Phi[\mathbf v]-\mathcal V_\Phi(\mathbf X,0)
 =\int_0^T(\mathbf v-\tfrac12H^{-1}P\mathbf x)^\top
 H(\mathbf v-\tfrac12H^{-1}P\mathbf x)\,dt\ge0,
\]
where $J_\Phi$ is the cost in Eq.~\eqref{eq:cost} at $t=0$.
This supplies a direct optimality check.

For exact liquidation, diagonalize
$H^{-1/2}CH^{-1/2}=U\operatorname{diag}(\mu_i)U^\top$ and set
$\kappa_i=\sqrt{\lambda\mu_i}$. The modes satisfy
$\ddot y_i=\kappa_i^2y_i$, giving
\begin{equation}
 \mathbf x(t)=H^{-1/2}U\operatorname{diag}\!\left(
 \frac{\sinh[\kappa_i(T-t)]}{\sinh(\kappa_iT)}\right)U^\top H^{1/2}\mathbf X.
 \label{eq:traj}
\end{equation}
The hard-limit feedback matrix is
$P(t)=2H^{1/2}U\operatorname{diag}(\kappa_i\coth[\kappa_i(T-t)])U^\top H^{1/2}$.
For zero-risk modes, use
$\sinh(\kappa s)/\sinh(\kappa T)\to s/T$ and
$\kappa\coth(\kappa s)\to1/s$.
Appendix~A of \cite{bergault2022} gives the $R=0$ Riccati solution;
Almgren--Chriss \cite{almgrenchriss2000}, App.~A.3, gives the classical
inventory modes.

For a finite penalty, a formula valid without covariance/penalty commutation is
\begin{equation}
 \begin{pmatrix}Y(s)\\Z(s)\end{pmatrix}
 =\exp\!\left[s\begin{pmatrix}0&\tfrac12H^{-1}\\2\lambda C&0\end{pmatrix}\right]
 \begin{pmatrix}I\\\Phi\end{pmatrix},\qquad P(T-s)=Z(s)Y(s)^{-1}.
 \label{eq:finite-penalty}
\end{equation}
Indeed, $Y'=H^{-1}Z/2$ and $Z'=2\lambda CY$ give
$(ZY^{-1})'=2\lambda C-(ZY^{-1})H^{-1}(ZY^{-1})/2$ and $P(T)=\Phi$.
This identity is derived here from Eq.~\eqref{eq:ric}. A scalar-mode
finite-penalty solution requires $H^{-1/2}\Phi H^{-1/2}$ and
$H^{-1/2}CH^{-1/2}$ to commute.

\subsection*{Scope and exit criterion}

$R=0$ removes mean reversion, not Brownian price noise. Price independence
uses the paper's terminal data and its $R=0$ solution. A feasible signed-control
optimum also solves the sales-only problem at that initial state; it does
not establish a sales-constrained value function over an entire inventory grid.

Step 2's source reduction is recorded. In Step 3, review its signs, factors
of two and terminal data before accepting the benchmark.

\section{Step 3: verify the classical benchmark (current)}
\label{sec:stage3}

\subsection*{Instance and derivation}

Use the synthetic, dimensionless instance
\[
 H=I_2,\qquad C=\begin{pmatrix}1&0.4\\0.4&1\end{pmatrix},\qquad
 \lambda=T=1,\qquad\mathbf X=(3,1)^\top,\qquad\mathbf x(1)=0.
\]
The orthonormal covariance eigenvectors are
$u_+=(1,1)^\top/\sqrt2$ and $u_-=(1,-1)^\top/\sqrt2$ with
\[
 \mu_+=1.4,\quad\mu_-=0.6,\qquad
 \kappa_+=\sqrt{1.4}=1.1832159566,\quad
 \kappa_-=\sqrt{0.6}=0.7745966692.
\]
Thus $C\succ0$, $\eta_j>0$, and both rates are real. Set $s=1-t$ and
\[
 f_\pm(s)=\frac{\sinh(\kappa_\pm s)}{\sinh\kappa_\pm},\qquad
 a_\pm(s)=\frac{\kappa_\pm\cosh(\kappa_\pm s)}{\sinh\kappa_\pm}.
\]
The trajectory and selling rates are
\begin{equation}
 x_1=2f_++f_-,\quad x_2=2f_+-f_-,\qquad
 v_1=2a_++a_-,\quad v_2=2a_+-a_-.
 \label{eq:stage3-rates}
\end{equation}
They give $\mathbf x(0)=(3,1)^\top$ and $\mathbf x(1)=0$.

For the equal-impact, equal-volatility family with $0<\rho<1$ and
$X_1=\theta X_2$, $\theta>1$, $X_2>0$, the potentially negative rate is
\[
 v_2=\frac{X_2}{2}\big((\theta+1)a_+-(\theta-1)a_-\big).
\]
The relative rate $\mathcal R(s)=a_-(s)/a_+(s)$ satisfies
\[
 \frac{d}{ds}\log \mathcal R(s)
 =\kappa_-\tanh(\kappa_-s)-\kappa_+\tanh(\kappa_+s)\le0,
\]
because $k\mapsto k\tanh(ks)$ increases for $k>0$, $s\ge0$.
Hence $\mathcal R(s)\le \mathcal R(0)$: checking the terminal-time ratio is sufficient and
necessary for this particular family, rather than an assumption about all
multi-asset problems. Specifically,
\begin{equation}
 \frac{\theta+1}{\theta-1}\ge
 \mathcal R(0)=\frac{\kappa_-\sinh\kappa_+}{\kappa_+\sinh\kappa_-}
 =1.133429601920,
 \qquad
 \theta\le\frac{\mathcal R(0)+1}{\mathcal R(0)-1}=15.989177597983.
 \label{eq:stage3-bound}
\end{equation}
The choice $\theta=3$ satisfies this bound. More explicitly, since
$a_+(s)\ge a_+(0)$ and $2-\mathcal R(s)\ge2-\mathcal R(0)>0$,
\[
 v_2(t)\ge a_+(0)(2-\mathcal R(0))=0.6931333103,
 \qquad v_1(t)\ge2a_+(0)+a_-(0)=2.5062988967.
\]
Both minima occur at $t=1$; these are one-sided terminal limits of the rates.
Thus both inventories stay nonnegative and decrease to zero for the entire
interval. A time-grid check alone would not establish this all-time claim.

\subsection*{Checks to reproduce}

Targets retained from the existing notes are
\[
 \mathbf v(0)=(4.0488436487,1.6639810858)^\top,\qquad
 J^*=\tfrac12\mathbf X^\top P(0)\mathbf X=13.810512031917.
\]
Recompute the cost by direct quadrature and compare it with the Riccati value.
The full numerical check script is absent; the table below remains a set
of targets to reproduce, not completed verification.

For an independent check, minimize
$\sum_{k=1}^N\|\mathbf x_{k-1}-\mathbf x_k\|^2/\tau
+\tau\sum_{k=1}^N\mathbf x_k^\top C\mathbf x_k$ by solving its block
tridiagonal stationarity system. Its Hessian is positive definite, so this
is the unique global discrete minimizer.
\begin{center}
\begin{tabular}{@{}rrr@{}}
\toprule
$N$ & Maximum inventory error vs. continuum & Discrete optimal cost \\
\midrule
16 & $6.945\times10^{-5}$ & 13.432845715492 \\
64 & $4.361\times10^{-6}$ & 13.714251771181 \\
256 & $2.726\times10^{-7}$ & 13.786331703645 \\
1024 & $1.704\times10^{-8}$ & 13.804459745777 \\
\bottomrule
\end{tabular}
\end{center}
The trajectory error decreases quadratically; the cost has first-order
right-endpoint quadrature error. These are different error measures.
The permanent-impact/spread offsets are not part of these reported costs.

\subsection*{Source, scope and exit criterion}

The trajectory comes from Step 2; the parameter choice and sales-only
inequality are derived here. The all-time proof applies to the stated
hard-liquidation instance. A finite penalty needs its own feasibility check,
and $H=I$ does not test a noncommuting implementation.

\textbf{Pending.} Review the Eq.~16 mapping, Riccati equation, terminal data
and sales-only proof. Reproduce the cost and convergence table, and compare
Eq.~\eqref{eq:finite-penalty} with direct Riccati integration for a
noncommuting penalty/covariance pair. Record the checks, then begin Step 4.

\section{Step 4: time-to-go and Cole--Hopf (next)}
\label{sec:stage4}

\subsection*{Derivation}

Set $s=T-t$. Add compatible inventory diffusion to the time-reversed HJB:
\begin{equation}
 W_s=\nu\nabla\!\cdot(H^{-1}\nabla W)
 -\tfrac14(\nabla W)^\top H^{-1}\nabla W+\lambda\mathbf x^\top C\mathbf x,
 \qquad W(\mathbf x,0)=\tfrac12\mathbf x^\top\Phi\mathbf x.
 \label{eq:viscous}
\end{equation}
For $\psi=e^{-W/(4\nu)}$,
\[
 \psi_s=-\frac{W_s}{4\nu}\psi,\qquad
 \nabla\!\cdot(H^{-1}\nabla\psi)
 =\psi\left[\frac{(\nabla W)^\top H^{-1}\nabla W}{16\nu^2}
 -\frac{\operatorname{tr}(H^{-1}\nabla^2W)}{4\nu}\right].
\]
The squared-gradient terms cancel, giving
\begin{equation}
 \boxed{\psi_s=-\hat H\psi,\qquad
 \hat H=-\nu\nabla\!\cdot(H^{-1}\nabla)
 +\frac{\lambda}{4\nu}\mathbf x^\top C\mathbf x,}
 \quad\psi(\mathbf x,0)=e^{-\mathbf x^\top\Phi\mathbf x/(8\nu)}.
 \label{eq:linearised}
\end{equation}
The control follows from $\nabla W=-4\nu\nabla\log\psi$:
\begin{equation}
 \boxed{\mathbf v^*=-2\nu H^{-1}\nabla\log\psi.}
 \label{eq:recovery}
\end{equation}
To check whether diffusion changes the desired control, substitute
$W=\tfrac12\mathbf x^\top P_s\mathbf x+c(s)$ into Eq.~\eqref{eq:viscous}:
\[
 P_s'=2\lambda C-\tfrac12P_sH^{-1}P_s,\qquad P_0=\Phi,
 \qquad c'(s)=\nu\operatorname{tr}(H^{-1}P_s),\quad c(0)=0.
\]
Thus $P_s=P(T-s)$ and $\nabla W=P_s\mathbf x$ agree with the original
finite-penalty control on the full space; diffusion changes only $c$.

\subsection*{Source, caveats and exit criterion}

Kappen \cite{kappen2005} supplies the exponential-transformation/noise-control
compatibility principle; the $H^{-1}$-weighted quadratic specialization is
derived here. This argument does not certify bounded-domain boundary
conditions or a globally sales-constrained HJB. Taking $\Phi\to\infty$
in the unnormalized heat initial data gives zero almost everywhere; use
a specified finite penalty or a properly normalized limiting kernel.

\textbf{To do after Step 3.} Review the cancellation and time orientation.
Choose finite $\Phi$ and $\nu$, solve the linear PDE classically, recover
Eq.~\eqref{eq:recovery}, and compare with the Riccati solution for that same
$\Phi$. For one asset the check is $-4\nu\partial_x\log\psi=P(t)x$.
Measure control error under grid refinement and domain enlargement.

\section{Step 5: spectrum and choosing the diffusion parameter}
\label{sec:stage5}

\subsection*{Derivation and source}

For $\lambda C\succ0$, the coordinates $\mathbf y=U^\top H^{1/2}\mathbf x$
from Step 2 make the operator a sum of harmonic oscillators:
\begin{equation}
 \hat H=\sum_i\left[-\nu\partial_{y_i}^2+
 \frac{\kappa_i^2}{4\nu}y_i^2\right],\qquad
 E_{\mathbf n}=\sum_i\kappa_i(n_i+\tfrac12),\qquad
 \ell_i=\sqrt{2\nu/\kappa_i}.
 \label{eq:oscillator}
\end{equation}
These follow from the standard full-space oscillator spectrum, with
$\kappa_i=\sqrt{\lambda\mu_i}$. The levels are decay rates in
$e^{-s\hat H}$; the ground-state contribution decays as $e^{-E_0s}$,
$E_0=\tfrac12\sum_i\kappa_i$.

For finite $\Phi$, the quadratic solution of Step 4 gives
\[
 \psi(\mathbf X,s)=\exp\!\left[-\frac{\mathbf X^\top P_s\mathbf X}{8\nu}
 -\frac14\int_0^s\operatorname{tr}(H^{-1}P_r)\,dr\right].
\]
Larger $\nu$ broadens the Gaussian. Its unnormalized value alone does not
predict readout probability, which depends on the normalized grid state.

\subsection*{Exit criterion}

Sweep $\nu$ and record $\ell_i$, $E_0T$, domain size, control error and
normalized probabilities near $\mathbf X$. Select $\nu$ using the grid and
readout costs together. The full-space spectrum does not certify a
truncated numerical domain.

\section{Step 6: discretisation}
\label{sec:stage6}

\subsection*{Construction and bound}

Use $N_x=2^{n_x}$ points per asset on $[-L,L)$, with $\Delta x=2L/N_x$.
Choose boundary conditions explicitly and enlarge $L$ to test truncation.
For a central second difference,
\[
 (D_{2,a}\psi)_{\mathbf j}
 =\frac{\psi_{\mathbf j+\mathbf e_a}-2\psi_{\mathbf j}
 +\psi_{\mathbf j-\mathbf e_a}}{\Delta x^2}.
\]
Let $\mathsf X_a$ be the diagonal matrix of asset-$a$ grid coordinates. Then
\begin{equation}
 \hat H_d=-\nu\sum_a\eta_a^{-1}D_{2,a}
 +\frac{\lambda}{4\nu}\sum_{a,b}C_{ab}\mathsf X_a\mathsf X_b,
 \label{eq:disc-operator}
\end{equation}
with the stencil/domain bound
\begin{equation}
 \|\hat H_d\|\lesssim\frac{4\nu}{\Delta x^2}\sum_a\eta_a^{-1}
 +\frac{\lambda L^2}{4\nu}\sum_{a,b}|C_{ab}|.
 \label{eq:norm-bound}
\end{equation}
The kinetic bound uses $\|-D_{2,a}\|\le4/\Delta x^2$; the potential bound
uses $|x_a|\le L$. These are standard finite-difference estimates applied
here, and must be checked for the chosen boundary stencil.

\subsection*{Exit criterion}

Preserve Hermiticity, compare the norm bound with numerical eigenvalues,
and quantify recovered-control error against Riccati as $\Delta x$, $L$
and $\nu$ vary. The final choice must resolve the Gaussian widths and
include the starting portfolio. No circuit-cost advantage is assumed from
the stencil alone.

\section{Step 7: Schr\"odingerisation}
\label{sec:stage7}

\subsection*{Construction and source}

The heat evolution $e^{-s\hat H_d}$ is nonunitary. Embed it in a larger
unitary system before applying quantum Hamiltonian simulation.
Following Jin, Liu and Yu \cite{jin2023}: with $A=-\hat H_d$ Hermitian
and negative semidefinite, set $w=e^{-p}\psi$ for an auxiliary variable $p$
for $p>0$, extend the initial data as
$w(0,p)=e^{-|p|}\psi(0)$ on the real line, and Fourier-transform in $p$
with frequencies $\eta_k$. This gives the
Schr\"odinger equation
\begin{equation}
  i\,\partial_s\hat w=H_S\hat w,\qquad H_S=-\hat H_d\otimes D_\eta,\qquad
  \|H_S\|=\|\hat H_d\|\,\eta_{\max},\ \ \eta_{\max}=\pi/\Delta p
  \label{eq:schrodinger}
\end{equation}
with $\psi=e^{p^*}w(\cdot,p^*)$ recovered at $p^*>0$.
The generator and extension follow Sec.~3.1, Eqs.~3.2--3.4 of
\cite{jin2023}. Derive the finite $p$-range, spacing and recovery probability
for this instance. A constant shift of $\hat H_d$ preserves
$\nabla\log\psi$, but does not remove the need to account for recovery cost.
\subsection*{Exit criterion}

Recover $\psi$ to the target accuracy in a small classical simulation of
the enlarged system, with auxiliary-domain error and success probability
included. The unitary generator alone does not close this step.

\section{Step 8: resource count}
\label{sec:stage8}

\subsection*{Formulas and sources}

Let $n_p=\log_2N_p$ encode the auxiliary coordinate, and let $n_{\rm anc}$
be the working qubits of the actual circuit. Then
\begin{equation}
 N_q=Mn_x+n_p+n_{\rm anc}.
 \label{eq:nq}
\end{equation}
For a block encoding with normalization $\alpha\ge\|H_S\|$, qubitization/QSVT
provides the asymptotic query bound \cite{lowchuang2019,gilyen2019}
\begin{equation}
 Q=O\!\left(\alpha T+\frac{\log(1/\varepsilon)}{\log\log(1/\varepsilon)}\right),
 \qquad N_T^{\rm simulation}=Q\,T_{\rm block}+T_{\rm synthesis}.
 \label{eq:queries}
\end{equation}
$T_{\rm block}$ is the T-cost of one block-encoding call and
$T_{\rm synthesis}$ counts synthesized phase rotations. The norm bound
in Step 6 does not itself construct a block encoding or determine its
normalization. State preparation, recovery and repeated readout calls must
be included in the total cost.

\subsection*{Exit criterion}

Specify the encoding, arithmetic precision and error allocation; derive
$\alpha$, $n_{\rm anc}$ and gate costs from that circuit. Produce a numerical
qubit/T-gate estimate for the benchmark and cross-check it at circuit level.
The query bound alone is not a numerical resource estimate.

\section{Step 9: readout}
\label{sec:stage9}

\subsection*{Observable and cost}

The output is the initial selling rate, recovered from nearby positive
amplitude ratios:
\[
 \partial_{x_j}\log\psi(\mathbf X)\approx
 \frac{\log\psi(\mathbf X+\Delta x\mathbf e_j)-\log\psi(\mathbf X)}{\Delta x}.
\]
This forward difference has first-order error, separate from the
second-order Laplacian stencil. Normalized amplitudes and probabilities are
\[
 a_j=\frac{\psi(\mathbf x_j)}{\sqrt{\sum_k|\psi(\mathbf x_k)|^2}},\qquad
 p_j=|a_j|^2.
\]
For a positive heat solution, probability ratios determine amplitude ratios.
Amplitude estimation \cite{brassard2002} resolves a small amplitude to
relative error $\epsilon_a$ with query scaling
$O(1/(\epsilon_a|a_j|))$, given coherent preparation access. Each call
includes preparation and recovery. Log-amplitude errors are divided by
$\Delta x$ and multiplied by $2\nu H^{-1}$ in the recovered control.

\subsection*{Exit criterion}

Allocate $\epsilon_a$ from the desired control error and include the resulting
readout calls in the total T-count. Report normalized probabilities near
$\mathbf X$; small probabilities can dominate the resource estimate.

\section{Step 10: validation and paper}
\label{sec:stage10}

\subsection*{Required results}

Validate the full Steps 4--9 route against the classical solution, including
the $\lambda\to0$ straight-line liquidation limit. Sweep
$\lambda$, correlation $\rho$, target accuracy $\varepsilon$, $\nu$ and $n_x$;
keep the asset count $M$ explicit. Record terminal penalty, boundary
conditions, numerical errors and all preparation/recovery/readout costs
beside the resource tables.

\subsection*{Writing notes and exit criterion}

Retain the source-to-notation map, the sales-only proof and the
finite-penalty comparison in the paper. Separate source results from the
specializations derived here. The benchmark is synthetic and classically
solvable; the paper's claim must match its resource estimates.

Before claiming novelty, compare the proposed route with the work combining
viscous Hamilton--Jacobi equations, Cole--Hopf and Schr\"odingerisation at
\url{https://arxiv.org/abs/2606.28752} and nonlinear-PDE simulation at
\url{https://arxiv.org/abs/2508.01640}. Detailed comparison remains open.

Completion requires a reproducible classical validation, an explicit circuit
and error budget, and resource tables that include readout. No downstream
numerical result is currently complete.

\begin{thebibliography}{9}
\bibitem{almgrenchriss2000} R.~Almgren and N.~Chriss,
\emph{Optimal Execution of Portfolio Transactions}, J.~Risk 3(2), 5--39
(2000), Appendix~A.3.

\bibitem{bergault2022} P.~Bergault, F.~Drissi and O.~Gu\'eant,
\emph{Multi-asset optimal execution and statistical arbitrage strategies
under Ornstein--Uhlenbeck dynamics}, SIAM J.~Financial Math.~13(1),
353--390 (2022). \url{https://arxiv.org/pdf/2103.13773v4}.

\bibitem{kappen2005} H.~J.~Kappen, \emph{Linear theory for control of
nonlinear stochastic systems}, Phys.~Rev.~Lett.~95, 200201 (2005).
\url{https://arxiv.org/abs/physics/0411119}.

\bibitem{jin2023} S.~Jin, N.~Liu and Y.~Yu,
\emph{Quantum simulation of partial differential equations via
Schr\"odingerisation: technical details}, Phys.~Rev.~A 108, 032603 (2023),
Sec.~3.1. \url{https://arxiv.org/abs/2212.14703}.

\bibitem{lowchuang2019} G.~H.~Low and I.~L.~Chuang,
\emph{Hamiltonian simulation by qubitization}, Quantum 3, 163 (2019).
\url{https://arxiv.org/abs/1610.06546}.

\bibitem{gilyen2019} A.~Gily\'en, Y.~Su, G.~H.~Low and N.~Wiebe,
\emph{Quantum singular value transformation and beyond}, STOC (2019).
\url{https://arxiv.org/abs/1806.01838}.

\bibitem{brassard2002} G.~Brassard, P.~H\o yer, M.~Mosca and A.~Tapp,
\emph{Quantum amplitude amplification and estimation}, Contemp.~Math.~305,
53--74 (2002). \url{https://arxiv.org/abs/quant-ph/0005055}.
\end{thebibliography}
\end{document}
